\section{Introduction}
\label{sec:related}
\IEEEPARstart{S}{atellite} synthetic aperture radar (SAR) supports land-cover monitoring through cloud cover and without daylight. Sentinel-1 observations provide complementary measurements of surface backscatter and interferometric coherence for this task~\cite{Jacob2020,Sica2019}. Using several orbital tracks expands the available observation schedule, but also changes how the radar illuminates the surface. Incidence angle, look direction, and terrain affect the measured backscatter; radiometric terrain correction addresses the associated variation in illuminated area~\cite{Small2011}. A classifier deployed across tracks must therefore use features that retain class information as the observation configuration changes. This requirement becomes particularly important when only a short sequence of acquisitions is available.

Repeat-pass interferometric coherence measures the magnitude of the normalized complex correlation between coregistered SAR acquisitions~\cite{Bamler1998,Touzi1999}. It adds sensitivity to changes in the scattering configuration over the interval between observations. Stable structures can retain correlation, whereas changes within vegetation reduce it. This contrast provides information complementary to backscatter power and has been used for forest mapping and the separation of forests from artificial surfaces~\cite{Sica2019,Zhuang2024,Aghababaei2026}. The two observables thus describe different aspects of the same scene: intensity records scattering strength, while coherence records the persistence of the complex return. Their combination offers a physical basis for improving class separation with a small number of acquisitions.

Much of the evidence for this complementarity comes from time-series classification. Jacob \textit{et al.}~\cite{Jacob2020} compared coherence-based representations across land-cover classifiers, while Mestre-Quereda \textit{et al.}~\cite{MestreQuereda2020} combined backscatter and coherence from a full year of Sentinel-1 observations for crop-type mapping. Recent studies connect interferometric observables to crop growth and use coherence time series to detect seeding dates~\cite{Yuan2024,ChenX2025}. Their value depends on when observations sample the vegetation cycle: L\"ow \textit{et al.}~\cite{Low2024} examined crop-specific signatures across years and orbits, directly connecting these observables to phenology.

Recent SAR classifiers exploit temporal structure with distribution-aware prototypes~\cite{Guo2025}, polarimetric feature sequences~\cite{LiuN2025}, and learned spatial--temporal representations~\cite{Russo2025,He2025}. Multisensor attention models combine these sequences with optical observations~\cite{Anandakrishnan2025}, while optical spatial--temporal fusion and parcel-based pipelines extend crop mapping to larger regions~\cite{Li2026,Sanchez2026}. Multifrequency experiments further show how polarimetric and interferometric measurements complement one another~\cite{Romero2026}. These approaches extract information along additional observation dimensions. The acquisition budget is therefore an essential part of any comparison of representation quality.

Scattering-based representations offer another route to class separation. Recent work uses dual-polarization target descriptors for Sentinel-1 clustering~\cite{Verma2025} and scattering powers for single-date crop classification with simulated NISAR data~\cite{Nihar2026}. Statistical approaches characterize class separability through intensity distributions~\cite{Negri2025} or geodesic polarimetric parameters~\cite{daSilva2026}. Multipass covariance modeling explicitly exploits temporal correlations when classifying scattering symmetries~\cite{Hanis2025}. Particularly relevant here, Zhuang \textit{et al.}~\cite{Zhuang2024} use repeat-pass coherence to distinguish vegetation volume scattering from rotated built-up structures. These results motivate testing temporal stability as a compact additional observable in dual-polarization land-cover mapping.

The accompanying intensity sequence also affects the measured benefit of coherence. In the 18-scene crop-mapping experiment of Zhao \textit{et al.}~\cite{Zhao2024}, dual-polarization backscatter achieved 95.12\% overall accuracy (OA), and adding VV coherence increased OA to 95.33\%. The resulting $0.21\pp$ increment describes the additional information available after a growing-season intensity sequence. A two-acquisition budget presents a different inference problem. It supplies one coherence measurement and two intensity observations, with no seasonal trajectory to learn. Moreover, those two intensity images still support spatial texture, polarization ratios, and gradients. Quantifying coherence's marginal value therefore requires a comparison against features extracted from the same image pair, including a developed intensity representation.

Cross-track deployment adds a second dimension to that comparison. Geographic transfer experiments have documented the dependence of crop-classification performance on the training and deployment regions~\cite{Orynbaikyzy2022}. Multisource domain generalization and frequency-based style mixing address distribution changes in remote-sensing image classification and segmentation~\cite{Han2024,Iizuka2024}; meta-learning also supports transfer between polarimetric SAR datasets and frequency bands~\cite{Luo2025}. These studies motivate separating the information supplied by an observable from the adaptation supplied by a learning algorithm. Here, observations from four tracks cover the same study area, allowing directional transfers across orbital configurations on common geographic support. Each track supplies its own repeat-pass pair for coherence estimation; the classifier is then transferred between tracks. Evaluating the corresponding in-track predictions gives a direct interpretation of the coherence increment. Gains in both settings indicate additional class discrimination, while a change in the difference between in-track and cross-track accuracy measures its effect on the transfer penalty.

This study quantifies the marginal classification value of a single 12-day coherence measurement under a fixed two-acquisition Sentinel-1 budget (Fig.~\ref{fig:workflow}). Four tracks yield twelve directional transfers, evaluated with random forest and support-vector classifiers. A five-representation design compares basic intensity features, a 28-dimensional intensity representation, and the addition of raw or spatially transformed coherence. All comparisons share the acquisition pairs and training budgets, with a 400\,m training-to-validation exclusion during spatial model selection. Five-class balanced accuracy measures the average class-level benefit, and paired spatial block bootstrap contrasts quantify its uncertainty. Combining these contrasts with in-track evaluation and class-level confusion analysis identifies where coherence contributes information and how that contribution affects cross-track deployment.

\begin{figure*}[!t]
\centering
\includegraphics[width=0.98\textwidth]{fig1_workflow.pdf}
\caption{Evaluation of coherence under a two-acquisition budget. (a) Shared SLC pairs for each track. (b) Five representations comparing intensity texture, raw coherence, and derived coherence features. (c) Spatial model selection with a 400\,m training-to-validation exclusion and additional 5\,km directional holdouts. (d) Classification and transfer contrasts with paired block resampling across compared representations.}
\label{fig:workflow}
\end{figure*}

